\section{Methodology}

\subsection{Pulsed Single-Photon LiDAR Model}

For a target at range $R$, the round-trip delay is

\begin{equation}
\tau = \frac{2R}{c},
\end{equation}

where $c$ is the speed of light. The target return is assigned to the corresponding time-of-flight histogram bin.

The transmitted photon budget during one sensing dwell is

\begin{equation}
N_{\mathrm{tx}} = N_{\mathrm{pulse}} N_{\mathrm{pulses}},
\end{equation}

where $N_{\mathrm{pulse}}$ denotes photons per pulse and $N_{\mathrm{pulses}}$ denotes the number of pulses accumulated during the dwell.

Two-way atmospheric transmission is modeled using the Beer--Lambert relationship

\begin{equation}
T^2(R) =
\exp\left(-2\alpha R\right),
\end{equation}

where $\alpha$ is the atmospheric extinction coefficient expressed in consistent distance units.

The beam radius at range $R$ is approximated by

\begin{equation}
r_b(R) =
r_0 + R\tan\left(\frac{\theta}{2}\right),
\end{equation}

where $r_0$ is the initial beam radius and $\theta$ is the full-angle divergence. The intercepted photon fraction is limited by the ratio between projected target area and beam-spot area.

The simulator additionally incorporates distributed atmospheric backscatter, background photons, detector efficiency, timing jitter, dark counts, dead time, pile-up, and afterpulsing.

\subsection{Sequential Target-Probability Update}

Each candidate region maintains a posterior target-presence probability. After observing $y$ photon counts in a range window, the posterior is updated using Poisson likelihoods under the target-absent hypothesis $H_0$ and target-present hypothesis $H_1$:

\begin{equation}
P(H_1 \mid y) =
\frac{
P(y \mid H_1)P(H_1)
}{
P(y \mid H_1)P(H_1)
+
P(y \mid H_0)P(H_0)
}.
\end{equation}

\subsection{Physics-Guided Information Allocation}

For region $i$, the allocation score is

\begin{equation}
S_i =
\mathcal{H}(p_i)
\log\left(1+\mathrm{SBR}_i\right)
\frac{1}{\sqrt{n_i+1}},
\end{equation}

where $\mathcal{H}(p_i)$ is the binary entropy of the target posterior, $\mathrm{SBR}_i$ is the predicted signal-to-background ratio, and $n_i$ is the number of previous measurements assigned to the region.

The next photon allocation is assigned to

\begin{equation}
i^{*} = \arg\max_i S_i.
\end{equation}

\subsection{Comparison Policies}

The proposed strategy is compared with:

\begin{enumerate}
    \item Uniform round-robin allocation.
    \item Random allocation.
    \item Greedy allocation based on maximum target posterior.
    \item Entropy allocation based on maximum posterior uncertainty.
\end{enumerate}

\subsection{Evaluation Protocol}

Six candidate regions are positioned between 500~m and 3000~m. Each policy is evaluated over five random seeds and six possible target locations, producing 30 episodes per policy. Each sensing action allocates $2.5\times10^{11}$ photons, with a maximum total budget of $5.0\times10^{12}$ photons. Detection occurs when the posterior probability of the true target region reaches 0.90.

\subsection{Simulation Parameters}

The principal transmitter, receiver, detector, histogram, and allocation
parameters are summarized in Table~\ref{tab:simulation_parameters}.

\begin{table}[ht]
\centering
\caption{Principal simulation and sensing parameters.}
\label{tab:simulation_parameters}
\begin{tabular}{lll}
\toprule
Category & Parameter & Value \\
\midrule
Laser & Wavelength & 1550 nm \\
Laser & Pulse width & 1 ns \\
Laser & Full-angle divergence & 0.5 mrad \\
Laser & Transmitter efficiency & 0.70 \\
Receiver & Aperture diameter & 0.10 m \\
Receiver & Optical efficiency & 0.55 \\
Detector & Photon-detection efficiency & 0.35 \\
Detector & Timing jitter & 80 ps \\
Detector & Dark-count rate & 500 Hz \\
Detector & Dead time & 50 ns \\
Detector & Afterpulse probability & 0.01 \\
Histogram & Number of bins & 4096 \\
Histogram & Bin width & 5 ns \\
Environment & Candidate ranges & 500--3000 m \\
Allocation & Photons per action & $2.5\times10^{11}$ \\
Allocation & Total photon budget & $5.0\times10^{12}$ \\
Decision & Detection threshold & 0.90 \\
Evaluation & Episodes per policy & 30 \\
\bottomrule
\end{tabular}
\end{table}

